%==============================================================
%  IEEE J-BHI submission — full draft
%  Compile on Overleaf with pdfLaTeX (IEEEtran is preinstalled).
%==============================================================
\documentclass[journal]{IEEEtran}

\usepackage{cite}
\usepackage{amsmath,amssymb,amsfonts}
\usepackage{algorithmic}
\usepackage{algorithm}
\usepackage{graphicx}
\usepackage{textcomp}
\usepackage{booktabs}
\usepackage{multirow}
\usepackage{url}
\hyphenation{fed-er-at-ed}

\begin{document}

\title{Absent Classes, Not Client Drift:\\
Diagnosing and Correcting Federated Collapse\\
in Brain Tumor MRI Classification}

\author{Zahraa Ibrahim Khaleel and Mohammed Sahib Mehdi%
\thanks{Manuscript received \today.}
\thanks{The authors are with the Department of Computer Science, College of Science, Al-Nahrain University, Baghdad, Iraq (e-mail: [YOUR EMAIL]).}}

\markboth{IEEE Journal of Biomedical and Health Informatics,~Vol.~XX, No.~X, 2026}%
{Khaleel \MakeLowercase{\textit{and}} Mehdi: Class-Presence-Guided Restricted Softmax}

\maketitle

\begin{abstract}
Federated learning allows several hospitals to train a shared diagnostic model
without exchanging patient images, but its accuracy degrades when the label
distribution differs across sites. In brain tumor classification this
degradation is severe, because a participating centre may never encounter one
of the tumor types at all. We show that under a Dirichlet partition with
concentration $0.3$, thirteen of seventeen simulated sites hold fewer than
three tumor types, and that a federated model trained with the standard
cross-entropy loss under these conditions loses six to seven points of balanced
accuracy and becomes acutely sensitive to how the partition happens to fall, its
worst partitions returning $0.447$ against a chance level of $0.333$. We trace the failure to a specific
term: for a class a client does not hold, the cross-entropy gradient on that
class logit is strictly positive at every step, so local training displaces the
corresponding output parameters in a direction fixed by the partition draw.
Removing that term entirely, by masking absent-class logits during local
updates, raises balanced accuracy from $0.584$ to $0.649$ over twenty seeds
(Holm-corrected $p=0.002$, Cohen's $d=0.91$), reduces the across-seed standard
deviation from $0.074$ to $0.028$, and closes $54\%$ of the gap to centralised
training on the pooled data. The advantage replicates on the BRISC~2025
benchmark, which uses independently annotated masks and a different class
balance. Comparison against logit calibration shows the two corrections to be
statistically indistinguishable ($p=0.12$), which supports the diagnosis rather
than any single remedy: what matters is that the absent-class term is removed,
not how. A controlled sweep confirms this, with the benefit scaling from $17.5$
points at high local absence to no measurable effect when every client holds
every class.
\end{abstract}

\begin{IEEEkeywords}
Federated learning, label distribution skew, restricted softmax, brain tumor,
magnetic resonance imaging, non-IID data, SCAFFOLD.
\end{IEEEkeywords}

%==============================================================
\section{Introduction}
\IEEEPARstart{M}{agnetic} resonance imaging is the primary modality for
characterising intracranial tumors, and automated classification of tumor type
from MRI has been studied extensively as a decision-support aid
\cite{cheng2015,gumaei2019,deepak2019}. Progress in this area depends on access
to large annotated cohorts, yet the images involved are among the most tightly
regulated categories of clinical data. Institutional review requirements and
national privacy legislation frequently prevent hospitals from pooling scans
into a central repository, which leaves individual centres training on datasets
too small to support a reliable model.

Federated learning offers a way around this constraint. Rather than moving
images, participating sites exchange model updates with a coordinating server,
so that patient data never leaves the institution where it was acquired
\cite{mcmahan2017,kairouz2021}. The approach has been evaluated for
multi-institutional brain tumor segmentation \cite{sheller2020} and has been
proposed as a general framework for collaborative clinical machine learning
\cite{rieke2020}.

The difficulty is that the classical convergence analysis of federated
averaging assumes that client datasets are drawn from a common distribution,
and clinical data violate this assumption in a specific and consequential way.
Referral patterns, departmental specialisation and scanner availability
determine which conditions a centre sees. A neurosurgical unit that treats
pituitary adenomas may accumulate hundreds of such cases and none of another
tumor type. The resulting label skew is not a mild statistical nuisance: when a
client holds no examples of class $c$, the gradient of the cross-entropy loss
with respect to the logit of $c$ is strictly positive at every optimisation
step, so local training drives that logit downward without any counteracting
signal. Averaging models that have each suppressed a different subset of
classes yields a global classifier whose final layer has been corrupted in a
manner that depends on how the partition happened to fall.

Several corrections operate directly on the softmax. FedRS rescales the logits
of absent classes by a factor $\alpha \in [0,1]$ during local training
\cite{fedrs2021}; FedLC generalises this to a per-class margin derived from
label frequency \cite{fedlc2022}; FedVLS adds a distillation term aimed
specifically at vacant classes \cite{fedvls2025}. These methods were developed
and tuned on natural-image benchmarks with large label spaces, where each client
is missing most of the classes. The regime encountered in tumor-type
classification is different: the label space contains three classes, and a
client typically misses one. Whether the tuned intermediate values of $\alpha$
reported for CIFAR-scale problems remain appropriate when the label space is
small has not been examined, and the behaviour of the limiting case
$\alpha = 0$ has not been characterised.

This paper makes the following contributions.

\begin{itemize}
\item We identify the mechanism behind federated collapse in this setting. The
cross-entropy gradient on a class the client does not hold is strictly positive
at every step, so local training displaces the corresponding output parameters
along a direction determined by the partition draw. We show that this term, and
not client drift in the shared representation, is what destabilises training,
and we quantify the condition that produces it: a Dirichlet partition with
concentration $0.3$ leaves $13$ of $17$ clients without at least one tumor type.

\item We show that removing the term is what matters, not the particular way it
is removed. Full masking of absent-class logits and frequency-based logit
calibration both recover $6$ to $7$ points of balanced accuracy over federated
averaging and are statistically indistinguishable from each other
($p=0.12$ over twenty seeds), while partial scaling at the setting published for
FedRS recovers less than one point. This locates the effect in the presence or
absence of the term rather than in any specific operator, and the comparison
against calibration is not reported in either originating paper.

\item We place the result on an absolute scale by measuring centralised
training on the pooled data. The correction closes $54\%$ of the gap between
federated averaging and the centralised upper bound, which no prior study of
this dataset reports.

\item We isolate the mechanism through a controlled sweep over partition
severity. The benefit scales monotonically with the mean fraction of classes
absent per client, from $17.5$ points at high absence to no measurable effect
when every client holds every class, which is the behaviour the gradient
argument predicts and a general regularisation account does not.

\item We report three negative findings that follow from the same analysis:
learning-rate warmup with cosine decay and extended training, both standard
remedies for unstable federated optimisation, increase the collapse rate under
this form of skew; the FedProx proximal term adds nothing to the corrected
objective; and the benefit halves when the backbone is enlarged from a three-block
convolutional network to a residual network, which bounds the regime in which
the correction is worth applying.
\end{itemize}

%==============================================================
\section{Related Work}

\subsection{Federated Optimisation under Heterogeneous Data}
Federated averaging aggregates locally trained weights in proportion to dataset
size \cite{mcmahan2017}. Its accuracy deteriorates when client distributions
diverge, an effect analysed empirically by Zhao \emph{et al.} \cite{zhao2018}
and confirmed across a range of partitions by Li \emph{et al.} \cite{li2022exp}.
Two families of correction dominate. FedProx constrains the local solution to
remain near the current global model through a proximal penalty
\cite{fedprox2020}. SCAFFOLD instead estimates the difference between local and
global gradient directions and subtracts it during local steps, which removes
client drift without restricting the size of the local update
\cite{scaffold2020}. Both target drift in the shared representation. Neither
addresses the case in which a client's local objective is defined over a
strict subset of the label space, which is the situation examined here.

Normalisation interacts with heterogeneity. Batch normalisation estimates
statistics from local batches, so under non-identical client distributions the
running estimates diverge and averaging becomes unreliable. Hsieh \emph{et al.}
showed that replacing batch normalisation with group normalisation
\cite{groupnorm2018} recovers most of the accuracy lost to this effect
\cite{hsieh2020}, a result reproduced in later normalisation benchmarks
\cite{casella2023}. FedBN takes the alternative route of keeping normalisation
layers local \cite{fedbn2021}. We adopt group normalisation for the reasons
given in \cite{hsieh2020}.

\subsection{Softmax and Logit Adjustment for Label Skew}
FedRS is the closest antecedent to the present work \cite{fedrs2021}. It
multiplies the logits of classes absent from a client by a scalar
$\alpha \in [0,1]$, which reduces but does not eliminate the gradient reaching
those logits; $\alpha = 1$ recovers federated averaging. The reported
experiments use CIFAR-10 and CIFAR-100 with $\alpha = 0.5$, and gains over
federated averaging range from $4.6$ to $14.9$ points depending on the partition
and the backbone. FedLC replaces the single scalar with a per-class margin
computed from local label frequencies, treating missing classes as the extreme
of the minority-class case, and reports larger gains on Dirichlet partitions of
CIFAR-10 \cite{fedlc2022}. FedVLS observes that calibration methods can fall
below federated averaging when the number of vacant classes is large, and adds a
distillation objective to compensate \cite{fedvls2025}.

\subsection{Relation to the Boundary Case Already Reported in FedRS}
\label{sec:fedrs_relation}
Because our correction is the boundary case of FedRS, we set out explicitly what
that paper already covers, so that the difference is a matter of record rather
than of emphasis.

Li and Zhan evaluate $\alpha \in \{0.0, 0.1, 0.5, 0.9, 1.0\}$, so the value we
use is among those they tried. They also treat it analytically, observing that
at $\alpha = 0$ the method ``degenerates into fixed proxies'' whose weights are
not updated and therefore do not contribute to the aggregation, which is the same
mechanism we describe in Section~\ref{sec:absentgrad}. Their empirical summary is
that $\alpha = 0$ obtains results comparable to $\alpha = 0.5$ in most cases and
can be worse with more local epochs. Consistent with that reading, $\alpha = 0$
appears in their convergence figures but not in their headline tables, and the
setting they carry forward is $\alpha = 0.5$. They further report a combined
variant, Scaffold/RS, which is the best configuration in their comparison,
so pairing the restricted objective with SCAFFOLD is likewise established rather
than new. Finally, in an appendix they observe weaker improvements with a VGG11
backbone than with a smaller network and conjecture that the restriction helps
more when the feature extractor is small; the backbone measurement in
Section~\ref{sec:backbone} is the controlled version of that conjecture.

Three things are therefore not claimed here: the operator, the pairing with
SCAFFOLD, and the qualitative idea that backbone size matters.

What differs is the regime and, with it, the answer. Their evaluation uses ten
to one hundred classes with two to twenty per client, where a client is missing
most of the label space and removing all of that negative evidence costs as much
as the protection gains; ``comparable to $\alpha = 0.5$'' is the expected outcome
there. Our setting has three classes with a mean local absence of $0.31$, so a
client typically lacks one class, the discarded evidence is small, and the
protection dominates. In that regime $\alpha = 0$ is not comparable to
$\alpha = 0.5$ but ahead of it by $5.8$ points
(Table~\ref{tab:main}). Their own caveat is consistent with this: the case they
found unfavourable is many local epochs, and we use one.

The contribution is accordingly not the boundary value but an account of what
governs it. Section~\ref{sec:sweep} measures the benefit as a function of local
class absence and finds it monotone, ranging from $17.5$ points to nothing, which
turns a reported equivalence and a measured advantage into two points on one
curve rather than a contradiction. To that we add a comparison against logit
calibration, which neither paper reports and which shows the two removals to be
indistinguishable; an absolute reference from centralised training; and
patient-level evaluation on medical images, where the missing-class condition
arises from referral patterns rather than from a synthetic partition.

\subsection{Brain Tumor Classification and Evaluation Practice}
The Figshare dataset introduced by Cheng \emph{et al.} \cite{cheng2015} is the
standard benchmark for three-class tumor-type classification and remains the
reference point for reported accuracy. Subsequent work has applied engineered
feature descriptors \cite{gumaei2019}, transfer learning from natural-image
networks \cite{deepak2019,swati2019} and capsule architectures
\cite{afshar2019}.

Comparison across these studies is complicated by evaluation protocol. Each
patient contributes multiple adjacent slices with near-identical appearance, so
a random slice-level split places highly correlated images in both training and
test sets. Yagis \emph{et al.} measured the consequence directly on brain MRI cohorts for
Alzheimer's and Parkinson's disease, and found that slice-level splitting
inflated reported accuracy by $29$ to $55$ points relative to subject-level
splitting on the same data, and that randomly labelled scans reached
approximately $96\%$ under a slice-level split against $50\%$ under a
subject-level split \cite{yagis2021}. The pathology differs from ours, but the
mechanism is a property of the data structure rather than of the disease: any
cohort in which one subject contributes many adjacent slices is exposed to it.
We therefore group by patient identifier throughout and treat published
slice-level figures as measuring a different quantity.

%==============================================================
\section{Methods}

\subsection{Problem Setting}
Let $K$ clients hold private datasets $\mathcal{D}_1,\dots,\mathcal{D}_K$ over a
common label space $\mathcal{Y}$ with $|\mathcal{Y}| = C$. Write
$\mathcal{P}_k \subseteq \mathcal{Y}$ for the set of classes that appear in
$\mathcal{D}_k$ and $\mathcal{A}_k = \mathcal{Y} \setminus \mathcal{P}_k$ for the
absent classes. Federated training minimises
\begin{equation}
\min_{\theta}\ \sum_{k=1}^{K} \frac{n_k}{n}\, F_k(\theta),
\qquad
F_k(\theta) = \frac{1}{n_k}\sum_{(x,y)\in\mathcal{D}_k} \ell(f_\theta(x), y),
\label{eq:objective}
\end{equation}
with $n_k = |\mathcal{D}_k|$ and $n = \sum_k n_k$. We define the \emph{local
absence} of client $k$ as $a_k = |\mathcal{A}_k| / C$ and report the mean
$\bar{a} = K^{-1}\sum_k a_k$ as the summary statistic of partition severity.

\subsection{The Absent-Class Gradient}
\label{sec:absentgrad}
Let $z = f_\theta(x) \in \mathbb{R}^C$ denote the logits and $p =
\mathrm{softmax}(z)$. For the cross-entropy loss $\ell(z,y) = -\log p_y$, the
derivative with respect to the logit of an arbitrary class $c$ is
\begin{equation}
\frac{\partial \ell}{\partial z_c} = p_c - \mathbb{1}[c = y].
\label{eq:grad}
\end{equation}
If $c \in \mathcal{A}_k$ then $\mathbb{1}[c=y] = 0$ for every sample in
$\mathcal{D}_k$, so the derivative reduces to $p_c > 0$ at every step. The
parameters of the output unit for class $c$ are therefore driven in one
direction throughout local training, and the magnitude of this displacement
depends on how many and which classes are absent, which is determined by the
partition and thus by the random seed. When the server averages models whose
output layers have been displaced along different class directions, the
resulting classifier is biased in a way that varies from run to run. This
accounts for both the loss of mean accuracy and the large across-seed variance
reported in Section~\ref{sec:results}.

\subsection{Class-Presence-Guided Restricted Softmax}
The correction rescales the contribution of absent classes to the local
partition function. For client $k$,
\begin{equation}
\ell_k(z,y) = -\log
\frac{\exp(z_y)}{\displaystyle\sum_{c \in \mathcal{Y}} s_c^{(k)} \exp(z_c)},
\qquad
s_c^{(k)} =
\begin{cases}
1, & c \in \mathcal{P}_k,\\[2pt]
\alpha, & c \in \mathcal{A}_k,
\end{cases}
\label{eq:rs}
\end{equation}
which recovers standard cross-entropy at $\alpha = 1$ and reproduces FedRS
\cite{fedrs2021} for intermediate $\alpha$. Differentiating \eqref{eq:rs} with
respect to an absent-class logit gives
\begin{equation}
\frac{\partial \ell_k}{\partial z_c}
= \frac{\alpha \exp(z_c)}{\sum_{c'} s_{c'}^{(k)}\exp(z_{c'})}
= \alpha\, \tilde{p}_c , \qquad c \in \mathcal{A}_k ,
\label{eq:grad_alpha}
\end{equation}
so the scaling factor acts directly on the rate at which the absent-class
parameters are displaced. Setting $\alpha = 0$ removes the term entirely: the
absent-class output parameters receive no gradient and are returned to the
server unchanged from the values at which they were broadcast. In our
implementation this is realised by adding $-\infty$ (numerically $-10^{9}$) to
the absent logits before the softmax, which is equivalent to \eqref{eq:rs} in the
limit and numerically stable.

Intermediate values of $\alpha$ slow the displacement without stopping it. In a
federated setting the effect accumulates, because each round re-broadcasts the
aggregated parameters and repeats the local pass. Over $T$ rounds the residual
displacement scales with $\alpha T$, so a factor that appears small in a single
epoch is not small after forty rounds. This is the reason we expect, and
observe, that $\alpha = 0$ dominates intermediate settings when local absence is
high.

The argument also predicts when the correction should make no difference. If
$\mathcal{A}_k = \emptyset$ for all $k$, the modified loss \eqref{eq:rs} is
identical to cross-entropy for any $\alpha$, and no effect of the correction
should be measurable. Section~\ref{sec:sweep} tests this prediction directly.

\subsection{Training Procedure}
Local optimisation uses SCAFFOLD \cite{scaffold2020}, whose control variates
correct drift in the shared feature extractor, combined with the restricted
objective \eqref{eq:rs}, which corrects the output layer. The two act on
different parts of the model and are complementary. All normalisation layers are
group normalisation \cite{groupnorm2018} following \cite{hsieh2020}. Local
sampling is class-balanced within each client, and gradients are clipped to unit
norm. Algorithm~\ref{alg:main} summarises the procedure.

\begin{algorithm}[t]
\caption{Class-Presence-Guided Federated Training}
\label{alg:main}
\begin{algorithmic}[1]
\STATE \textbf{Input:} clients $1..K$, rounds $T$, learning rate $\eta$, scaling $\alpha$
\STATE Server initialises $\theta^{0}$ and control variate $c^{0}$
\FOR{$t = 0$ \TO $T-1$}
  \STATE Server broadcasts $\theta^{t}$, $c^{t}$
  \FOR{each client $k$ \textbf{in parallel}}
    \STATE $\mathcal{P}_k \leftarrow$ set of classes present in $\mathcal{D}_k$
    \STATE $\theta_k \leftarrow \theta^{t}$
    \FOR{each local batch $(x,y)$}
      \STATE $g \leftarrow \nabla_{\theta_k}\, \ell_k(f_{\theta_k}(x), y)$ using \eqref{eq:rs}
      \STATE $g \leftarrow \mathrm{clip}\big(g + c^{t} - c_k,\ 1\big)$
      \STATE $\theta_k \leftarrow \theta_k - \eta g$
    \ENDFOR
    \STATE Update $c_k$; return $\Delta_k = \theta_k - \theta^{t}$
  \ENDFOR
  \STATE $\theta^{t+1} \leftarrow \theta^{t} + \sum_k (n_k/n)\,\Delta_k$
  \STATE $c^{t+1} \leftarrow c^{t} + K^{-1}\sum_k (c_k - c^{t})$
\ENDFOR
\end{algorithmic}
\end{algorithm}

\subsection{Evaluation Protocol}
\label{sec:protocol}
Test patients are held out before partitioning, and all slices belonging to a
patient are assigned to the same side of the split, so no patient contributes
to both training and evaluation \cite{yagis2021}.

The primary metric is balanced accuracy, the unweighted mean of per-class
recall. The class distribution of the Figshare dataset is uneven
($708$ meningioma, $1426$ glioma, $930$ pituitary), so a classifier that
predicts the majority class for every input attains $46.5\%$ plain accuracy
while contributing nothing diagnostically; the same classifier scores $0.333$
under balanced accuracy. Since the failure mode under study is precisely
collapse onto a single class, plain accuracy would obscure the phenomenon the
experiments are designed to measure. We define collapse as a run with balanced
accuracy at or below $0.40$.

Every configuration is run over $20$ seeds, with the seed controlling both the
Dirichlet draw and model initialisation, so that comparisons are paired.
Differences are assessed with the Wilcoxon signed-rank test, which does not
assume normality of the per-seed differences, and reported with Cohen's $d$ for
paired samples and a $10{,}000$-resample bootstrap confidence interval. Because
each method is compared against three baselines, $p$-values are adjusted with
the Holm step-down procedure \cite{holm1979}.

%==============================================================
\section{Experimental Setup}

\subsection{Datasets}
The primary dataset is the Figshare collection of Cheng \emph{et al.}
\cite{cheng2015}, comprising $3064$ contrast-enhanced T1-weighted slices from
$233$ patients, distributed as $708$ meningioma, $1426$ glioma and $930$
pituitary, with expert tumor masks and patient identifiers.

External validation uses BRISC \cite{brisc2025}, released with expert
segmentation masks for the same three tumor types. We use $3000$ slices,
$1000$ per class, drawn from the segmentation subset. BRISC does not distribute
patient identifiers, so clients are formed by drawing slices rather than
patients; this limitation is stated again in Section~\ref{sec:limits}. We also
note that BRISC is assembled in part from sources that include the Cheng
collection, so the two benchmarks are not fully independent. The masks are newly
annotated and the class balance differs, which makes the replication informative
as a robustness check rather than as evidence of external generalisation in the
strict sense.

\subsection{Partitioning and Model}
Training patients are distributed over $K = 18$ clients by drawing, for each
class, a vector of proportions from a Dirichlet distribution with concentration
$\alpha_{\mathrm{dir}}$. Clients receiving fewer than four samples are removed,
which leaves $17$ active clients at $\alpha_{\mathrm{dir}} = 0.3$. At this
setting $13$ of $17$ clients hold fewer than three tumor types, and $\bar{a} =
0.31$. Fig.~\ref{fig:problem} shows the resulting distribution of absent classes
and the composition of each client.

\begin{figure*}[t]
\centering
\includegraphics[width=0.86\textwidth]{fig_problem.png}
\caption{The missing-class condition at $\alpha_{\mathrm{dir}}=0.3$. (a) Only
four of seventeen clients hold all three tumor types; nine are missing one and
four are missing two. (b) Class composition of each client, where empty cells
indicate a tumor type the client never observes during training.}
\label{fig:problem}
\end{figure*}

The classifier is a three-block convolutional network with $16$, $32$ and $64$
channels, group normalisation after each convolution, global average pooling and
a linear output layer. Images are resized to $128 \times 128$ and scaled to
$[0,1]$. Local optimisation uses stochastic gradient descent with learning rate
$0.01$, one local epoch per round and a batch size of $32$; training runs for
$40$ rounds. The deliberately small backbone keeps the comparison focused on the
loss function rather than on representational capacity. Section~\ref{sec:backbone}
measures what changes when the backbone is enlarged.

Two notation conventions are used throughout: $\alpha_{\mathrm{dir}}$ denotes
the Dirichlet concentration governing partition severity, and $\alpha$ without
subscript denotes the restricted-softmax scaling factor of \eqref{eq:rs}.

\subsection{Reproducibility}
Every configuration is executed from a fixed seed that determines the Dirichlet
draw, the train and test assignment of patients, the model initialisation and
the order of local batches. Methods compared within a table therefore see
identical partitions, which is what makes the paired tests in
Section~\ref{sec:protocol} valid. Results are written incrementally to disk so
that interrupted runs resume without repeating completed cells, and the
aggregate tables are regenerated from those records rather than transcribed.
Statistical tests are computed from the per-seed values, not from the summary
statistics.

We note one methodological check that shaped the design of the evaluation. An
earlier version of this study assessed per-patient adaptation using a support
and query split of each held-out patient's slices. Because a patient carries a
single tumor type across all slices, the support set determines the label of the
query set, and a predictor that ignores the images entirely and copies the
majority support label attains a balanced accuracy of $1.000$. Per-patient
adaptation is therefore not measurable on this task, and we report no such
figures.

%==============================================================
\section{Results}
\label{sec:results}

\subsection{Which Interventions Stabilise Training}
We first isolate the effective component by enabling one intervention at a time
on top of an uncorrected federated baseline (Table~\ref{tab:ablation} and
Fig.~\ref{fig:ablation}). This ablation was run at an earlier and shorter
schedule than the comparisons that follow, twenty rounds with two local epochs
rather than forty rounds with one, so its baseline sits at $0.551$ rather than
the $0.584$ reported in Table~\ref{tab:main}. The two numbers are not in
conflict; they describe the same objective under different budgets, and only
within-table comparisons are made. Gradient
clipping alone does not help. Learning-rate warmup with cosine decay and
extended training, both standard responses to unstable federated optimisation,
increase the collapse rate from zero to $15$ and $14$ runs out of $20$
respectively. Class-balanced local sampling leaves the mean essentially
unchanged. The restricted softmax raises balanced accuracy from $0.551$ to
$0.611$ and eliminates collapse. Adding the FedProx proximal term on top of the
corrected objective changes neither the mean nor the variance.

\begin{table}[t]
\centering
\caption{Component ablation on Figshare over $20$ seeds. Collapse denotes a run
with balanced accuracy $\leq 0.40$.}
\label{tab:ablation}
\begin{tabular}{lcc}
\toprule
Configuration & Balanced acc. & Collapses \\
\midrule
Federated baseline            & $0.551 \pm 0.031$ & $0/20$ \\
\quad + gradient clipping     & $0.529 \pm 0.052$ & $1/20$ \\
\quad + warmup, cosine decay  & $0.371 \pm 0.065$ & $15/20$ \\
\quad + extended training     & $0.379 \pm 0.068$ & $14/20$ \\
\quad + balanced sampling     & $0.544 \pm 0.081$ & $1/20$ \\
\quad + restricted softmax    & $\mathbf{0.611 \pm 0.055}$ & $0/20$ \\
\quad + FedProx term          & $0.611 \pm 0.055$ & $0/20$ \\
\bottomrule
\end{tabular}
\end{table}

\begin{figure}[t]
\centering
\includegraphics[width=0.88\columnwidth]{fig_ablation.png}
\caption{Component ablation. Only the restricted softmax (bottom) improves on
the uncorrected baseline. Learning-rate warmup with cosine decay and extended
training push the mean toward the chance level of $0.333$ (dotted line).}
\label{fig:ablation}
\end{figure}

The two configurations that fail share a property: both increase the number of
local optimisation steps taken against the uncorrected objective, either by
lengthening the schedule or by raising the effective step size in the later
rounds. Under \eqref{eq:grad}, more steps against that objective means more
displacement of the absent-class parameters, so the direction of these results
follows from the mechanism rather than contradicting it.

\subsection{Comparison against Softmax Corrections}
Table~\ref{tab:main} compares four local objectives on identical partitions.
Federated averaging and FedRS at its published setting $\alpha=0.5$ behave
almost alike, separated by less than one point. The two objectives that remove
the absent-class term outright, full masking and frequency-based logit
calibration, both gain six to seven points over federated averaging, and the
difference between them is not significant ($\Delta=0.004$, Holm $p=0.108$,
masking ahead on $8$ of the $20$ partitions and calibration on the other $12$).

We read this pattern as evidence for the diagnosis rather than for a particular
operator. Partial scaling attenuates the absent-class gradient by a factor
$\alpha$ but leaves it in the objective, and over forty rounds the residual
displacement accumulates; the result is close to no correction at all. Once the
term is removed, whether by masking or by a margin large enough to be equivalent
to masking, the method used to remove it makes no measurable difference. The
practical distinction between the two is that masking requires no
hyperparameter, whereas calibration introduces a temperature that must be
tuned; in a federated clinical deployment, where no held-out validation set
spans all sites, this is a meaningful simplification.

The reduction in across-seed standard deviation is as informative as the change
in the mean. Federated averaging and FedRS vary with standard deviation $0.074$
and $0.072$, while the two corrected objectives sit at $0.028$ and $0.026$. The
worst partition under the hard mask returns $0.598$, above the \emph{mean} of
either uncorrected baseline.

Trained on the pooled data, the same backbone reaches $0.705$. The correction
therefore closes $54\%$ of the distance between federated averaging and
centralised training, which locates the remaining gap in the federation itself
rather than in the loss function.

\begin{table}[t]
\centering
\caption{Figshare, patient-level balanced accuracy over $20$ seeds. Differences
are computed against the hard mask; $p$-values are Holm-adjusted across the
three comparisons. The centralised bound is trained on the pooled data with the
same backbone and step budget ($5$ seeds).}
\label{tab:main}
\begin{tabular}{lcccc}
\toprule
Method & Balanced acc. & $\Delta$ & $p$ & $d$ \\
\midrule
FedAvg                  & $0.584 \pm 0.074$ & $-0.065$ & $0.002$ & $0.91$ \\
FedRS, $\alpha=0.5$     & $0.592 \pm 0.072$ & $-0.058$ & $0.002$ & $0.86$ \\
FedLC                   & $0.653 \pm 0.026$ & $+0.004$ & $0.123$ & $0.34$ \\
Hard mask, $\alpha=0$   & $\mathbf{0.649 \pm 0.028}$ & --- & --- & --- \\
\midrule
\textit{Centralised bound} & \textit{$0.705 \pm 0.011$} & --- & --- & --- \\
\bottomrule
\end{tabular}
\end{table}

\subsection{Distribution of Outcomes across Partitions}
\begin{figure*}[t]
\centering
\includegraphics[width=0.86\textwidth]{fig_perseed.png}
\caption{Paired per-seed comparison against FedRS at $\alpha=0.5$. Each line
connects the two objectives on one partition; the heavy line joins the means.
The partitions on which the uncorrected baseline falls below $0.55$ are those
with the highest local absence, and these are the partitions on which removing
the absent-class term has the largest effect.}
\label{fig:perseed}
\end{figure*}

\begin{figure}[t]
\centering
\includegraphics[width=0.88\columnwidth]{fig_variance.png}
\caption{Standard deviation of balanced accuracy across seeds. Objectives that
remove the absent-class term give the narrowest spread on both benchmarks.}
\label{fig:variance}
\end{figure}

The aggregate figures in Table~\ref{tab:main} conceal a difference in the shape
of the distributions. Federated averaging and FedRS produce a bimodal pattern:
on most partitions they reach the mid-$0.6$ range, but on a minority they fall
to the high $0.4$s. Seeds $5$, $13$ and $14$ illustrate this, returning $0.449$,
$0.488$ and $0.447$ under federated averaging and $0.494$, $0.487$ and $0.447$
under FedRS. On the same three partitions the hard mask returns $0.689$, $0.636$
and $0.670$, and calibration returns $0.682$, $0.632$ and $0.661$. The two
corrected objectives track each other closely on precisely the partitions where
the uncorrected ones fail, which is the pattern the diagnosis predicts. The
lowest result over all twenty seeds under the hard mask is $0.598$, above the
mean of either uncorrected baseline. Fig.~\ref{fig:perseed} shows the paired
outcomes seed by seed, and Fig.~\ref{fig:variance} summarises the resulting
spread on both benchmarks.

This pattern is what the mechanism predicts. The seeds on which the baselines
fail are those whose Dirichlet draw concentrates entire classes on a small
number of clients, so that the remaining clients carry a large absent-class
gradient. Since the hard mask removes that gradient regardless of how the draw
falls, its outcome depends much less on the partition. Reporting only the mean
would understate the difference between the methods, because the methods differ
most in the tail.

The hard mask wins on $15$ of $20$ paired comparisons against FedRS on Figshare.
The five partitions on which it does not win are all cases where every method
performs well, and the differences there are within $0.02$. No partition was
found on which the hard mask fails while a baseline succeeds.

\subsection{Replication on BRISC}
Table~\ref{tab:brisc} repeats the comparison on BRISC. The ordering of the four
methods is preserved and the margin over FedRS at $\alpha = 0.5$ is $5.2$
points, close to the $5.8$ observed on Figshare. No method collapses on this
benchmark, which is consistent with its balanced class distribution, and the
hard mask nevertheless retains its advantage. The correction therefore does more
than prevent degenerate solutions.

\begin{table}[t]
\centering
\caption{BRISC, balanced accuracy over $20$ seeds. Differences are computed
against the hard mask, Holm-adjusted.}
\label{tab:brisc}
\begin{tabular}{lcccc}
\toprule
Method & Balanced acc. & $\Delta$ & $p$ & $d$ \\
\midrule
FedAvg                  & $0.621 \pm 0.061$ & $-0.055$ & $0.0015$ & $0.87$ \\
FedRS, $\alpha=0.5$     & $0.624 \pm 0.062$ & $-0.052$ & $0.0015$ & $0.82$ \\
FedRS, $\alpha=0.25$    & $0.635 \pm 0.061$ & $-0.041$ & $0.0018$ & $0.77$ \\
Hard mask, $\alpha=0$   & $\mathbf{0.676 \pm 0.046}$ & --- & --- & --- \\
\bottomrule
\end{tabular}
\end{table}

\begin{figure*}[t]
\centering
\includegraphics[width=0.86\textwidth]{fig_compare.png}
\caption{Balanced accuracy over $20$ seeds on both benchmarks. Error bars show
one standard deviation across seeds. The ordering of methods is preserved
between datasets, and the hard mask combines the highest mean with the smallest
spread in each case.}
\label{fig:compare}
\end{figure*}

\subsection{The Margin Scales with Local Absence}
\label{sec:sweep}
Section~III argued that the benefit of $\alpha = 0$ should depend on how many
classes are absent from each client and should disappear when none are. We test
this by varying $\alpha_{\mathrm{dir}}$ over $\{0.1, 0.5, 3.0\}$, which produces
mean local absence $\bar{a} \in \{0.50, 0.31, 0.07\}$, and evaluating
$\alpha \in \{0, 0.5, 1\}$ at each level over eight seeds
(Table~\ref{tab:sweep}).

\begin{table}[t]
\centering
\caption{Balanced accuracy as a function of partition severity and scaling
factor, eight seeds per cell. The final column reports the margin of
$\alpha = 0$ over $\alpha = 0.5$ with its paired Wilcoxon $p$-value.}
\label{tab:sweep}
\begin{tabular}{ccccccc}
\toprule
$\alpha_{\mathrm{dir}}$ & $\bar{a}$ & $\alpha{=}0$ & $\alpha{=}0.5$ & $\alpha{=}1$ & Margin & $p$ \\
\midrule
$0.1$ & $0.50$ & $0.629$ & $0.454$ & $0.471$ & $+0.175$ & $0.008$ \\
$0.5$ & $0.31$ & $0.670$ & $0.607$ & $0.602$ & $+0.063$ & $0.023$ \\
$3.0$ & $0.07$ & $0.659$ & $0.657$ & $0.656$ & $+0.002$ & $0.836$ \\
\bottomrule
\end{tabular}
\end{table}

\begin{figure}[t]
\centering
\includegraphics[width=0.85\columnwidth]{fig_absence.png}
\caption{Margin of $\alpha=0$ over $\alpha=0.5$ as a function of mean local
class absence. Error bars show the standard error of the paired difference
across eight seeds. The effect disappears when clients hold nearly the complete
label space.}
\label{fig:absence}
\end{figure}

The margin decreases monotonically with local absence and is no longer
distinguishable from zero when clients hold nearly the full label space. We
emphasise that $\alpha = 0$ remains the best of the three settings at every
level tested; what changes is the size of the advantage, not its sign. The
correction is therefore inactive rather than harmful in the balanced regime,
which is the behaviour \eqref{eq:rs} predicts when $\mathcal{A}_k = \emptyset$.

\subsection{Sensitivity to Backbone Capacity}
\label{sec:backbone}
The corrections studied here act on the output layer, so their effect should
depend on how much of the classification is carried by that layer. We repeated
the comparison with a residual network of roughly ten times the parameter count,
keeping the partitions, budget and normalisation unchanged
(Table~\ref{tab:backbone}).

\begin{table}[t]
\centering
\caption{Effect of backbone capacity, $10$ seeds. Gains are measured against
federated averaging on the same backbone.}
\label{tab:backbone}
\begin{tabular}{lcccc}
\toprule
& \multicolumn{2}{c}{Three-block CNN} & \multicolumn{2}{c}{Residual network}\\
\cmidrule(lr){2-3}\cmidrule(lr){4-5}
Method & Bal. acc. & Gain & Bal. acc. & Gain\\
\midrule
FedAvg      & $0.584$ & --- & $0.559$ & ---\\
FedRS $0.5$ & $0.592$ & $+0.8$ & $0.563$ & $+0.3$\\
FedLC       & $0.653$ & $+6.9$ & $0.593$ & $+3.4$\\
Hard mask   & $0.649$ & $+6.5$ & $0.588$ & $+2.9$\\
\bottomrule
\end{tabular}
\end{table}

Two observations follow. The benefit of removing the absent-class term halves
when the backbone is enlarged, from $6.5$ to $2.9$ points for the hard mask and
from $6.9$ to $3.4$ for calibration. This is the behaviour expected of a
correction that acts on the output layer: as more of the classification is
carried by the feature extractor, the share attributable to the final layer
falls, and with it the effect of altering that layer's objective. The
measurement bounds the regime in which our result should be expected to hold.

The residual network is also worse in absolute terms on every objective. With
roughly $2400$ training slices divided among seventeen clients, the larger model
has too little data per client to justify its capacity, and the ordering of the
methods is preserved but compressed. We therefore report the three-block network
as the primary configuration, a choice the comparison supports rather than
assumes.

%==============================================================
\section{Discussion}

The experiments support a claim about a mechanism rather than about a method.
Federated training fails here because clients optimise an objective defined over
classes they do not hold, and the size of the failure tracks how much of the
label space is missing locally. Three results converge on this reading.

Partial and full removal of the absent-class term behave very differently.
Scaling the term by $\alpha=0.5$ recovers less than one point, while removing it
recovers six to seven. If the benefit came from regularisation in a general
sense, a factor of two in the coefficient would not produce a factor of seven in
the outcome. What separates the two is that attenuation leaves a residual
gradient which accumulates over rounds, whereas removal does not.

Two unrelated operators that both remove the term perform the same. Masking sets
the absent-class contribution to zero directly; calibration subtracts a margin
large enough to have the same effect. They differ in form, in the parameters
they require, and in the literature they come from, yet their outcomes cannot be
distinguished at twenty seeds. This is what one expects if the operative
variable is the presence of the term rather than the identity of the correction,
and it is why we frame the contribution as a diagnosis. A practitioner may
reasonably choose either; masking is simply the one that needs no tuning.

The benefit disappears exactly when the mechanism says it should. At
$\bar{a}=0.07$, where almost every client holds every class, the corrected and
uncorrected objectives are within $0.2$ points and the difference is not
significant. A method that helped through some general smoothing effect would
not switch off so cleanly.

The variance reduction deserves separate comment. The displacement of the
absent-class parameters depends on which classes each client happens to lack,
which is a property of the partition draw. Removing the term removes a
seed-dependent error source, and the standard deviation falls from $0.074$ to
$0.028$ while the mean rises by six points. For deployment this is arguably the
more relevant quantity: a federation is assembled once and its participants are
not resampled, so what matters is the accuracy obtained on the particular
partition at hand, and the correction narrows the range of outcomes that
partition can produce.

The negative results are consistent with the same account. Warmup with cosine
decay and extended training both increase the number of local steps taken
against an objective that still contains the absent-class term, which
\eqref{eq:grad} identifies as the harmful quantity. That these are standard
remedies for federated instability indicates that the failure mode here is
distinct from the client drift they were designed for. FedProx points the same
way: a proximal penalty limits how far the local solution travels but not the
direction in which it travels, so it cannot prevent a displacement that is
systematic rather than excessive.

\section{Limitations}
\label{sec:limits}

Several constraints bound these conclusions.

The correction we recommend is not novel as an operator, and neither is its
pairing with SCAFFOLD. Both appear in the paper that introduced the scaling
rule, which evaluates the boundary value among its settings and reports a
combined Scaffold/RS variant; logit calibration was proposed separately as a
generalisation of the same idea. Our claim concerns which property of these
objectives does the work, and the regime in which the boundary case is the right
choice rather than merely an equivalent one. Readers seeking a new operator will
not find one here. Section~\ref{sec:fedrs_relation} states the overlap in full.

Both datasets consist of contrast-enhanced T1-weighted slices, so the results do
not speak to multi-sequence or volumetric protocols. The two benchmarks are also
not fully independent: BRISC is assembled from sources that include the Cheng
collection. Its masks are independently annotated and its class balance differs,
which makes the replication a robustness check rather than validation on an
unseen population.

BRISC does not distribute patient identifiers, so clients on that benchmark are
formed by drawing slices. Slices from one patient may be split across clients and
across the train and test partitions, which is the leakage condition avoided on
Figshare. The BRISC figures are therefore comparable across methods, since every
method faces the same partition, but not comparable in absolute terms to the
patient-level Figshare figures.

The benefit halves with a larger backbone, as Table~\ref{tab:backbone} shows. Our
results are established for models of the size appropriate to the data available
per client in this setting; they should not be extrapolated to federations with
enough data per site to support a large feature extractor.

Balanced accuracy weights the three tumor types equally, which is appropriate for
measuring discriminative capacity but does not reflect the asymmetric cost of
diagnostic errors in clinical use. A cost-sensitive evaluation would require a
clinically validated cost matrix, which we did not have access to.

Finally, we tested one alternative account and rejected it. An auxiliary
contrastive term with targets derived from continuous fractal descriptors,
intended to organise the latent space so that absent classes retain structure,
gave a marginal and unstable gain when added to the mask and performed ten points
worse when used in its place, with collapse on two partitions. The failure of a
representation-space remedy is further evidence that the pathology resides in the
output layer.

\section{Conclusion}

Federated brain tumor classification collapses under severe label skew because
clients optimise a loss defined over classes they do not hold, which displaces
the corresponding output parameters in a direction fixed by the partition draw.
Removing that term raises balanced accuracy by six to seven points over
federated averaging, cuts across-seed variance by more than half, and closes
$54\%$ of the gap to centralised training. Two different operators that remove
it perform equivalently, while attenuating it recovers almost nothing, which
places the effect in the term itself rather than in any particular remedy. The
benefit scales with the fraction of the label space absent at each client and
vanishes when no classes are missing. For federations of the size encountered in
tumor-type classification, masking absent-class logits achieves this without
introducing a hyperparameter.

\begin{thebibliography}{00}

\bibitem{cheng2015}
J. Cheng, W. Huang, S. Cao, R. Yang, W. Yang, Z. Yun, Z. Wang, and Q. Feng,
``Enhanced performance of brain tumor classification via tumor region
augmentation and partition,'' \emph{PLoS ONE}, vol.~10, no.~10, art.~e0140381,
2015, doi: 10.1371/journal.pone.0140381.

\bibitem{gumaei2019}
A. Gumaei, M. M. Hassan, M. R. Hassan, A. Alelaiwi, and G. Fortino, ``A hybrid
feature extraction method with regularized extreme learning machine for brain
tumor classification,'' \emph{IEEE Access}, vol.~7, pp.~36266--36273, 2019.

\bibitem{deepak2019}
S. Deepak and P. M. Ameer, ``Brain tumor classification using deep CNN features
via transfer learning,'' \emph{Comput. Biol. Med.}, vol.~111, p.~103345, 2019.

\bibitem{mcmahan2017}
B. McMahan, E. Moore, D. Ramage, S. Hampson, and B. A. y Arcas,
``Communication-efficient learning of deep networks from decentralized data,''
in \emph{Proc. Int. Conf. Artif. Intell. Statist. (AISTATS)}, 2017,
pp.~1273--1282.

\bibitem{kairouz2021}
P. Kairouz \emph{et al.}, ``Advances and open problems in federated learning,''
\emph{Found. Trends Mach. Learn.}, vol.~14, no.~1--2, pp.~1--210, 2021.

\bibitem{sheller2020}
M. J. Sheller, B. Edwards, G. A. Reina, J. Martin, S. Pati, A. Kotrotsou,
M. Milchenko, W. Xu, D. Marcus, R. R. Colen, and S. Bakas, ``Federated learning
in medicine: facilitating multi-institutional collaborations without sharing
patient data,'' \emph{Sci. Rep.}, vol.~10, art.~12598, 2020.

\bibitem{rieke2020}
N. Rieke \emph{et al.}, ``The future of digital health with federated
learning,'' \emph{npj Digit. Med.}, vol.~3, art.~119, 2020.

\bibitem{fedrs2021}
X.-C. Li and D.-C. Zhan, ``FedRS: Federated learning with restricted softmax for
label distribution non-IID data,'' in \emph{Proc. 27th ACM SIGKDD Conf. Knowl.
Discov. Data Min. (KDD '21)}, Virtual Event, Singapore, Aug.~2021,
pp.~995--1005, doi: 10.1145/3447548.3467254.

\bibitem{fedlc2022}
J. Zhang, Z. Li, B. Li, J. Xu, S. Wu, S. Ding, and C. Wu, ``Federated learning
with label distribution skew via logits calibration,'' in \emph{Proc. Int. Conf.
Mach. Learn. (ICML)}, 2022, pp.~26311--26329.

\bibitem{fedvls2025}
K. Guo, Y. Ding, J. Liang, Z. Wang, R. He, and T. Tan, ``Exploring vacant
classes in label-skewed federated learning,'' in \emph{Proc. AAAI Conf. Artif.
Intell.}, vol.~39, 2025, pp.~16960--16968.

\bibitem{zhao2018}
Y. Zhao, M. Li, L. Lai, N. Suda, D. Civin, and V. Chandra, ``Federated learning
with non-IID data,'' arXiv preprint arXiv:1806.00582, 2018.

\bibitem{li2022exp}
Q. Li, Y. Diao, Q. Chen, and B. He, ``Federated learning on non-IID data silos:
an experimental study,'' in \emph{Proc. IEEE Int. Conf. Data Eng. (ICDE)}, 2022,
pp.~965--978.

\bibitem{fedprox2020}
T. Li, A. K. Sahu, M. Zaheer, M. Sanjabi, A. Talwalkar, and V. Smith,
``Federated optimization in heterogeneous networks,'' in \emph{Proc. Mach.
Learn. Syst. (MLSys)}, 2020, pp.~429--450.

\bibitem{scaffold2020}
S. P. Karimireddy, S. Kale, M. Mohri, S. J. Reddi, S. U. Stich, and A. T. Suresh,
``SCAFFOLD: Stochastic controlled averaging for federated learning,'' in
\emph{Proc. Int. Conf. Mach. Learn. (ICML)}, 2020, pp.~5132--5143.

\bibitem{groupnorm2018}
Y. Wu and K. He, ``Group normalization,'' in \emph{Computer Vision -- ECCV
2018}, LNCS vol.~11217, 2018, pp.~3--19.

\bibitem{hsieh2020}
K. Hsieh, A. Phanishayee, O. Mutlu, and P. B. Gibbons, ``The non-IID data
quagmire of decentralized machine learning,'' in \emph{Proc. Int. Conf. Mach.
Learn. (ICML)}, 2020, pp.~4387--4398.

\bibitem{casella2023}
B. Casella, R. Esposito, A. Sciarappa, C. Cavazzoni, and M. Aldinucci,
``Experimenting with normalization layers in federated learning on non-IID
scenarios,'' \emph{IEEE Access}, vol.~12, pp.~47961--47971, 2024.

\bibitem{fedbn2021}
X. Li, M. Jiang, X. Zhang, M. Kamp, and Q. Dou, ``FedBN: Federated learning on
non-IID features via local batch normalization,'' in \emph{Proc. Int. Conf.
Learn. Represent. (ICLR)}, 2021.

\bibitem{swati2019}
Z. N. K. Swati \emph{et al.}, ``Brain tumor classification for MR images using
transfer learning and fine-tuning,'' \emph{Comput. Med. Imaging Graph.},
vol.~75, pp.~34--46, 2019.

\bibitem{afshar2019}
P. Afshar, K. N. Plataniotis, and A. Mohammadi, ``Capsule networks for brain
tumor classification based on MRI images and coarse tumor boundaries,'' in
\emph{Proc. IEEE Int. Conf. Acoust. Speech Signal Process. (ICASSP)}, 2019,
pp.~1368--1372.

\bibitem{yagis2021}
E. Yagis \emph{et al.}, ``Effect of data leakage in brain MRI classification
using 2D convolutional neural networks,'' \emph{Sci. Rep.}, vol.~11, art.~22544,
2021.

\bibitem{brisc2025}
A. Fateh, Y. Rezvani, S. Moayedi, S. Rezvani, F. Fateh, M. Fateh, and
V. Abolghasemi, ``BRISC: Annotated dataset for brain tumor segmentation and
classification,'' \emph{Sci. Data}, vol.~13, art.~361, Feb.~2026,
doi: 10.1038/s41597-026-06753-y.

\bibitem{holm1979}
S. Holm, ``A simple sequentially rejective multiple test procedure,''
\emph{Scand. J. Statist.}, vol.~6, no.~2, pp.~65--70, 1979.

\end{thebibliography}

\end{document}
